% !PRINTED 651
\weberpage{651}
12. \gesperrt{Jedes ganze Funktional $\bar\omega$ ist der größte gemeinschaftliche
 Teiler zweier Funktionen $\alpha, \beta$ in $\Omega$.}

Um dies zu beweisen, nehmen wir $\alpha = \bar\omega \mu$, teilbar durch $\bar\omega$,
dann $\beta$ teilbar durch $\bar\omega$ und $\beta : \bar\omega$ relativ prim zu $\mu$; dann ist
\[
\bar\omega = \alpha x + \beta y
\]
der größte gemeinschaftliche Teiler von $\alpha$ und $\beta$.

\begin{center}
\textbf{§ 183.\qquad Basen und Basisformen der Funktionale.}
\end{center}

Es sei jetzt
\weq{1}{\omega_1, \omega_2, \ldots, \omega_n}
eine Minimalbasis des Körpers $\Omega$ (§ 176) und
\weq{2}{\alpha_1, \alpha_2, \ldots, \alpha_n}
eine andere Basis von $\Omega$. Die lineare Substitution, durch die
die $\alpha$ mit dem $\omega$ zusammenhängen, sei:
\weq{3}{\alpha_s = \overset{i}{\Sigma} a_{s,i} \omega_i,}
worin die $a_{s,i}$ ganze Funktionen in $Z$ sind, deren Determinante
\weq{4}{A = \Sigma \pm a_{1,1} a_{2,2} \ldots a_{n,n}}
nicht identisch verschwindet. Sind $t_1, t_2, \ldots, t_n$ Variable, so ist
\weq{5}{\lambda = \alpha_1 t_1 + \alpha_2 t_2 + \cdots + \alpha_n t_n}
der größte gemeinschaftliche Teiler von $\alpha_1, \alpha_2, \ldots, \alpha_n$.

Setzt man für $t_1, t_2, \ldots, t_n$ ganze Funktionen in $Z$, so entstehen
 aus $\lambda$ ganze Funktionen in $\Omega$, die alle durch $\lambda$ teilbar sind.

1. \gesperrt{Man nennt $\lambda$ eine Basisform und $\alpha_1, \alpha_2, \ldots, \alpha_n$
eine Basis des Funktionals $\lambda$, wenn man alle durch
$\lambda$ teilbaren Zahlen in $\Omega$ dadurch erhält, daß man
für $t_i$ ganze Funktionen in $Z$ setzt.}

Bilden die $\alpha_i$ eine solche Basis, so kann man ganze Funktionen,
$g_{r,s}^{(i)}$ in $Z$, so bestimmen, daß
\weq{6}{\alpha_r \omega_s = \overset{i}{\Sigma} g_{r,s}^{(i)} \alpha_i}
wird.

Daraus folgt:
\weq{7}{\lambda \omega_s = \overset{i}{\Sigma} \alpha_i t_{i,s},}
wenn
\weq{8}{t_{i,s} = \overset{r}{\Sigma} g_{r,s}^{(i)} t_r}
%CONT>>
